% Mathématiques 5e — cours V3
% Proportionnalité et pourcentages

\section{Objectifs}

\begin{itemize}
\item Utiliser des proportions et des pourcentages.
\item Calculer et appliquer une proportion ou un pourcentage.
\item Identifier une situation de proportionnalité dans un contexte concret.
\item Utiliser un coefficient de proportionnalité.
\item Résoudre par linéarité multiplicative, linéarité additive ou retour à l'unité.
\item Représenter une situation de proportionnalité par un tableau ou un graphique.
\item Reconnaître une situation proportionnelle à partir d'un tableau ou d'un graphique.
\item Reconnaître graphiquement un nuage de points associé ou non à une situation de proportionnalité.
\end{itemize}

\section{1. Deux grandeurs proportionnelles}

Deux grandeurs sont proportionnelles lorsque l'on passe de l'une à l'autre en multipliant toujours par le même nombre.

Ce nombre est le coefficient de proportionnalité.

\begin{definition}
Si une grandeur $y$ est proportionnelle à une grandeur $x$, il existe un nombre $k$ tel que :
\[
y=kx.
\]
\end{definition}

En 5e, cette écriture sert à formaliser une relation déjà rencontrée dans des situations concrètes.

\section{2. Contexte de prix}

Un kilogramme de tomates coûte :
\[
3{,}20\ \text{€}.
\]

Pour $2\ \text{kg}$ :
\[
2\times3{,}20=6{,}40\ \text{€}.
\]

Pour $4{,}5\ \text{kg}$ :
\[
4{,}5\times3{,}20=14{,}40\ \text{€}.
\]

Le coefficient de proportionnalité est :
\[
3{,}20.
\]

\section{3. Retour à l'unité}

Si $5$ objets coûtent :
\[
17{,}50\ \text{€},
\]
le prix d'un objet est :
\[
17{,}50\div5=3{,}50\ \text{€}.
\]

Pour $8$ objets :
\[
8\times3{,}50=28\ \text{€}.
\]

Le retour à l'unité reste une méthode fondamentale.

\section{4. Linéarité multiplicative}

Si l'on double une valeur de la première grandeur, on double la valeur correspondante de la seconde.

Par exemple :
\[
3\ \text{kg}\longrightarrow7{,}50\ \text{€}.
\]

Alors :
\[
6\ \text{kg}\longrightarrow15\ \text{€}.
\]

On a multiplié les deux grandeurs par :
\[
2.
\]

\section{5. Linéarité additive}

On peut additionner des correspondances.

Si :
\[
2\ \text{kg}\longrightarrow5\ \text{€},
\]
et :
\[
3\ \text{kg}\longrightarrow7{,}50\ \text{€},
\]
alors :
\[
5\ \text{kg}\longrightarrow12{,}50\ \text{€}.
\]

Cette propriété peut simplifier certains problèmes.

\section{6. Coefficient de proportionnalité}

Le coefficient de proportionnalité permet de passer directement d'une grandeur à l'autre.

\begin{exemple}
Une voiture parcourt :
\[
90\ \text{km}
\]
en :
\[
1\ \text{h}
\]
à vitesse moyenne constante.

Le coefficient reliant la durée en heures à la distance est :
\[
90.
\]

Pour :
\[
2{,}5\ \text{h},
\]
on obtient :
\[
90\times2{,}5=225\ \text{km}.
\]
\end{exemple}

\section{7. Tableau de proportionnalité}

On peut organiser les valeurs dans un tableau.

\begin{tabular}{c|c|c|c}
Masse en kg & $1$ & $3$ & $4{,}5$ \\
Prix en € & $3{,}20$ & $9{,}60$ & $14{,}40$ \\
\end{tabular}

Le même coefficient :
\[
3{,}20
\]
permet de passer de la première ligne à la seconde.

\section{8. Reconnaître un tableau proportionnel}

Un tableau est proportionnel si les quotients correspondants sont égaux.

Par exemple :
\[
\dfrac{6}{2}=3,
\]
\[
\dfrac{15}{5}=3,
\]
\[
\dfrac{24}{8}=3.
\]

Le coefficient est constant.

\begin{attention}
Un tableau à deux lignes n'est pas automatiquement un tableau de proportionnalité.
\end{attention}

\section{9. Exemple non proportionnel}

Une location coûte :
\[
4\ \text{€}
\]
de frais fixes puis :
\[
2\ \text{€}
\]
par heure.

Pour $1$ heure :
\[
6\ \text{€}.
\]

Pour $2$ heures :
\[
8\ \text{€}.
\]

Si la situation était proportionnelle, doubler la durée doublerait le prix :
\[
2\times6=12.
\]

Or :
\[
8\neq12.
\]

La situation n'est donc pas proportionnelle.

\section{10. Proportion}

Une proportion compare l'effectif d'une partie à l'effectif total.

\begin{definition}
Si une catégorie contient $n_A$ éléments parmi $n$ éléments au total, sa proportion est :
\[
\boxed{p=\dfrac{n_A}{n}}.
\]
\end{definition}

Par exemple, $9$ élèves sur $30$ :
\[
p=\dfrac9{30}=0{,}3.
\]

\section{11. Pourcentage}

Un pourcentage est une proportion exprimée sur $100$.

\coursefigure{/images/mathematiques/5eme/proportion-pourcentage.svg}{Grille de cent cases dont vingt-cinq sont coloriées.}{Vingt-cinq pour cent correspond à vingt-cinq parts sur cent, donc à un quart.}

On a :
\[
25\%=\dfrac{25}{100}=0{,}25.
\]

\section{12. Calculer un pourcentage}

Pour calculer :
\[
p\%
\]
d'une quantité $Q$, on peut utiliser :
\[
\boxed{\dfrac{p}{100}\times Q}.
\]

\begin{exemple}
Calculons :
\[
30\%
\]
de :
\[
80.
\]

On obtient :
\[
\dfrac{30}{100}\times80.
\]

Donc :
\[
0{,}3\times80=24.
\]

Ainsi :
\[
\boxed{30\%\text{ de }80=24}.
\]
\end{exemple}

\section{13. Retrouver un pourcentage}

Dans une classe de :
\[
24
\]
élèves, $6$ ont choisi une option.

La proportion vaut :
\[
\dfrac6{24}=\dfrac14.
\]

Donc :
\[
\dfrac14=25\%.
\]

Ainsi :
\[
\boxed{25\%}
\]
des élèves ont choisi cette option.

\section{14. Élection}

Quatre candidats obtiennent respectivement :
\[
6,\quad12,\quad3,\quad3
\]
voix.

Le total vaut :
\[
24.
\]

Pour le candidat ayant $12$ voix :
\[
\dfrac{12}{24}=\dfrac12=50\%.
\]

Pour un candidat ayant $3$ voix :
\[
\dfrac3{24}=\dfrac18=12{,}5\%.
\]

\section{15. Échelle}

Une échelle relie une longueur représentée à une longueur réelle.

À l'échelle :
\[
1:50\,000,
\]
un centimètre sur la carte correspond à :
\[
50\,000\ \text{cm}.
\]

Or :
\[
50\,000\ \text{cm}=500\ \text{m}.
\]

Donc :
\[
1\ \text{cm}\longrightarrow500\ \text{m}.
\]

\section{16. Exemple d'échelle}

Sur une carte au :
\[
1:25\,000,
\]
deux points sont séparés de :
\[
8\ \text{cm}.
\]

La distance réelle vaut :
\[
8\times25\,000=200\,000\ \text{cm}.
\]

Donc :
\[
200\,000\ \text{cm}=2\,000\ \text{m}=2\ \text{km}.
\]

Ainsi :
\[
\boxed{2\ \text{km}}.
\]

\section{17. Graphique de proportionnalité}

Une situation de proportionnalité peut être représentée par des points dans un repère.

\coursefigure{/images/mathematiques/5eme/proportion-graphique.svg}{Deux nuages de points, l'un aligné avec l'origine et l'autre non.}{Dans une situation proportionnelle, les points sont alignés avec l'origine du repère.}

\begin{propriete}
Un graphique représentant une situation de proportionnalité est constitué de points alignés avec l'origine.
\end{propriete}

\section{18. Reconnaître graphiquement}

Si les points d'un nuage sont alignés mais que la droite ne passe pas par l'origine, la situation n'est pas proportionnelle.

Si les points ne sont pas alignés, elle ne l'est pas non plus.

L'origine :
\[
O(0;0)
\]
est donc essentielle.

\section{19. Données discrètes}

En 5e, on peut rencontrer seulement quelques couples de valeurs.

Par exemple :
\[
(1;4),\quad(2;8),\quad(3;12),\quad(5;20).
\]

Tous ces points sont alignés avec l'origine.

Le coefficient est :
\[
4.
\]

On reconnaît donc une proportionnalité.

\section{20. Exemple développé}

Une recette utilise :
\[
350\ \text{g}
\]
de farine pour :
\[
4
\]
personnes.

Pour :
\[
10
\]
personnes, on peut revenir à l'unité :
\[
350\div4=87{,}5\ \text{g}.
\]

Puis :
\[
87{,}5\times10=875\ \text{g}.
\]

Donc :
\[
\boxed{875\ \text{g}}.
\]

\section{21. Choisir une procédure}

\begin{methode}
Dans un problème de proportionnalité :
\begin{enumerate}
\item identifier les deux grandeurs ;
\item vérifier que le contexte est proportionnel ;
\item choisir une procédure : multiplication, addition de colonnes, retour à l'unité ou coefficient ;
\item effectuer le calcul ;
\item contrôler l'unité et la vraisemblance.
\end{enumerate}
\end{methode}

\section{22. Tableau ou graphique}

Le tableau est pratique pour :
\begin{itemize}
\item calculer ;
\item organiser des couples de valeurs ;
\item repérer un coefficient.
\end{itemize}

Le graphique permet :
\begin{itemize}
\item de visualiser l'alignement ;
\item de comparer rapidement plusieurs situations ;
\item de reconnaître le passage par l'origine.
\end{itemize}

\section{23. Erreurs fréquentes}

\begin{itemize}
\item Supposer qu'une situation est proportionnelle sans vérification.
\item Confondre proportion et effectif.
\item Calculer $20\%$ en multipliant par $20$ au lieu de $0{,}20$.
\item Utiliser un coefficient différent selon les colonnes.
\item Oublier les unités dans une échelle.
\item Conclure à la proportionnalité parce que les points sont seulement alignés.
\item Oublier que la droite doit passer par l'origine.
\end{itemize}

\section{À retenir}

\begin{aretenir}
Deux grandeurs sont proportionnelles si un même coefficient multiplicatif les relie.

On peut résoudre avec :
\[
\boxed{\text{linéarité, retour à l'unité ou coefficient}}.
\]

Une proportion vaut :
\[
\boxed{
p=\dfrac{\text{partie}}{\text{total}}
}.
\]

Un pourcentage s'écrit :
\[
\boxed{
p\%=\dfrac{p}{100}
}.
\]

Graphiquement, une situation proportionnelle correspond à des points :
\[
\boxed{\text{alignés avec l'origine}}.
\]
\end{aretenir}
